\pagebreak
\section*{Comparison and Branch Condition Encoding}
\addcontentsline{toc}{section}{Comparison and Branch Condition Encoding}

This appendix explains the condition encoding and implementation philosophy used by the \texttt{CMP}, \texttt{CBR}, \texttt{ABR}, and \texttt{MBR} instructions. The design goal is to minimize hardware complexity by reusing a single subtraction/comparator path. Comparison operators are encoded in a compact and symmetric way that makes condition inversion trivial.

The following instructions use the same 3-bit comparison field (\texttt{cond}):

\begin{itemize}
  \item \textbf{CMP} — Compares two operands using the 3-bit \texttt{cond} field and produces a boolean result (0/1) in a register.
  \item \textbf{CBR} — Compares two registers and branches based on the same 3-bit \texttt{cond} field. This performs a direct comparison without requiring a preceding \texttt{CMP}.
  \item \textbf{ABR} — Performs an add operation and branches using a 3-bit \texttt{cond} field that tests the \emph{add result} against zero or other simple predicates. Intended for loop counters or pointer walking.
  \item \textbf{MBR} — Performs a move operation and branches using the same 3-bit \texttt{cond} field. The comparison is applied to the source operand, typically against zero.
\end{itemize}

All four instructions share a uniform 3-bit condition encoding, allowing common decoder and comparator logic.  
Branch target computation uses a sign-extended 15-bit offset shifted left by two, allowing word-aligned branch destinations within a \texttt{±64\,KB} range relative to the current instruction address.

\paragraph{Design Rationale}

The condition encoding is designed for symmetry and minimal logic cost. Conditions are arranged in complementary pairs so that flipping a single bit yields the inverse relation. This symmetry makes branch inversion trivial in hardware and simplifies assembler implementation.

\begin{center}
\begin{tabular}{|c|c|c|l|}
\hline
\textbf{Value} & \textbf{Binary} & \textbf{Mnemonic} & \textbf{Meaning / Branch if} \\ \hline
0 & 000 & \texttt{EQ} & Equal (\texttt{==}) \\ \hline
1 & 001 & \texttt{LT} & Less than (\texttt{<}) \\ \hline
2 & 010 & \texttt{GT} & Greater than (\texttt{>}) \\ \hline
3 & 011 & \texttt{EV} & Even (bit~0 = 0) \\ \hline
4 & 100 & \texttt{NE} & Not equal (\texttt{!=}) \\ \hline
5 & 101 & \texttt{GE} & Greater than or equal (\texttt{>=}) \\ \hline
6 & 110 & \texttt{LE} & Less than or equal (\texttt{<=}) \\ \hline
7 & 111 & \texttt{OD} & Odd (bit~0 = 1) \\ \hline
\end{tabular}
\end{center}

The most significant bit (\texttt{cond[2]}) acts as an \emph{invert bit}: flipping it inverts the sense of the comparison. The relationships between base and inverted conditions are as follows:

\begin{center}
\begin{tabular}{|c|c|c|}
\hline
\textbf{Base Condition} & \textbf{Inverted Condition} & \textbf{Encoding Pair} \\ \hline
\texttt{EQ} & \texttt{NE} & 000 / 100 \\ \hline
\texttt{LT} & \texttt{GE} & 001 / 101 \\ \hline
\texttt{GT} & \texttt{LE} & 010 / 110 \\ \hline
\texttt{EV} & \texttt{OD} & 011 / 111 \\ \hline
\end{tabular}
\end{center}

Since Twin-64 consistently works on signed d64-but quantities, all comparisons are signed by default. The \texttt{EV}/\texttt{OD} conditions allow efficient branching on bit~0, which is useful for address alignment tests, distinguishing even/odd indices, and loop unrolling.
